\subsection{Depth and property \texorpdfstring{$(S_k)$.}{(Sk)}}
\phantomsection
\label{subsection:IV.5.7}

\oldpage[IV-2]{103}

\begin{definition}[5.7.1]
\phantomsection
\label{IV.5.7.1}
Let $X$ be a locally Noetherian prescheme and let $\sh F$ be a
coherent $\sh O_X$-Module.  The depth \textup{(}respectively,
codepth\textup{)} of $\sh F$ at a point $x\in X$ is the number
$\operatorname{prof}(\sh F_x)$ \textup{(}respectively,
$\operatorname{coprof}(\sh F_x)$\textup{)}
\sref[0]{0.16.4.5}, \sref[0]{0.16.4.9}.  The codepth of
$\sh F$ is the number
\begin{equation}
\operatorname{coprof}(\sh F)
 =\sup_{x\in X}\operatorname{coprof}(\sh F_x).
\tag{5.7.1.1}\label{IV.5.7.1.1}
\end{equation}

We say that $\sh F$ is a Cohen--Macaulay $\sh O_X$-Module at the point
$x$ if $\sh F_x$ is a Cohen--Macaulay $\sh O_x$-module, that is
\sref[0]{0.16.5.1}, if
$\operatorname{coprof}(\sh F_x)=0$.  We say that $\sh F$ is a
Cohen--Macaulay $\sh O_X$-Module if it is so at every point,
in other words if $\operatorname{coprof}(\sh F)=0$.  A point
$x\in X$ such that $\sh O_x$ is a Cohen--Macaulay ring is also
called a Cohen--Macaulay point of $X$.
\end{definition}

The \emph{codepth} of $X$, denoted by $\operatorname{coprof}(X)$,
is the number $\operatorname{coprof}(\sh O_X)$.  We say that $X$
is a \emph{Cohen--Macaulay prescheme} if $\sh O_X$ is a
Cohen--Macaulay $\sh O_X$-Module, in other words if
$\operatorname{coprof}(X)=0$.  Every locally Noetherian prescheme of
\emph{dimension} $0$ is evidently a Cohen--Macaulay prescheme.  To say that
$\Spec(A)$ is a Cohen--Macaulay scheme means that $A$ is a Cohen--Macaulay
ring \sref[0]{0.16.5.13}.

\begin{definition}[5.7.2]
\phantomsection
\label{IV.5.7.2}
Let $X$ be a locally Noetherian prescheme, let $\sh F$ be a coherent
$\sh O_X$-Module, and let $k$ be an integer, positive or negative.
We say that $\sh F$ has property $(S_k)$ if, for every
$x\in X$, we have
\begin{equation}
\operatorname{prof}(\sh F_x)
 \geq \inf(k,\dim(\sh F_x)).
\tag{5.7.2.1}\label{IV.5.7.2.1}
\end{equation}
We say that $\sh F$ has property $(S_k)$ at a point $x\in X$
if, for every generization $x'$ of $x$ in $X$, we have
\begin{equation}
\operatorname{prof}(\sh F_{x'})
 \geq \inf(k,\dim(\sh F_{x'})).
\tag{5.7.2.2}\label{IV.5.7.2.2}
\end{equation}
\end{definition}

We say that $X$ satisfies property $(S_k)$
(resp.\ satisfies property $(S_k)$ at a point $x$) if $\sh O_X$ satisfies
property $(S_k)$ (resp.\ satisfies property $(S_k)$ at the point $x$).

For $\sh F$ to satisfy property $(S_k)$, it is evidently necessary and
sufficient that it satisfy it at every point of $X$.  If $U$ is an open
subset of $X$ and if $\sh F$ satisfies $(S_k)$, so does $\sh F|U$;
conversely, if $(U_\alpha)$ is an open covering of $X$ and if
$\sh F|U_\alpha$ satisfies $(S_k)$ for every $\alpha$, then
$\sh F$ satisfies $(S_k)$.

\begin{remarks}[5.7.3]
\phantomsection
\label{IV.5.7.3}
\begin{enumerate}[{\rm(i)}]
\item
\phantomsection
\label{IV.5.7.3.i}
Recall that we always have
$\operatorname{prof}(\sh F_x)\leq\dim(\sh F_x)$
\sref[0]{0.16.4.5.1} if $\sh F_x\neq0$.  To say that $\sh F$ has
property $(S_k)$ therefore means that we have
$\operatorname{prof}(\sh F_x)\geq k$ \emph{except} at the points $x\in X$
such that $\dim(\sh F_x)<k$, and that at these latter points we have
$\dim(\sh F_x)=\operatorname{prof}(\sh F_x)$, that is
\sref[0]{0.16.5.1}, that $\sh F$ is a \emph{Cohen--Macaulay}
$\sh O_X$-Module at these points; note that at the points where
$\dim(\sh F_x)=k$, we have $\operatorname{prof}(\sh F_x)=k$,
\oldpage[IV-2]{104}
so $\sh F$ is again a Cohen--Macaulay $\sh O_X$-Module at these
points.  To say that $\sh F$ has property $(S_k)$ \emph{for every}
$k$ therefore means that $\sh F$ is a \emph{Cohen--Macaulay
$\sh O_X$-Module}.  It is clear that for $k'\geq k$, property
$(S_{k'})$ implies $(S_k)$; for $k\leq0$, \emph{every} coherent
$\sh O_X$-Module has property $(S_k)$.

\item
\phantomsection
\label{IV.5.7.3.ii}
To verify condition \eqref{IV.5.7.2.1}, we may restrict ourselves to
the case where $x\in\operatorname{Supp}(\sh F)$; in the contrary case we
indeed have $\dim(\sh F_x)=-\infty$ \sref[0]{0.14.1.2}.

\item
\phantomsection
\label{IV.5.7.3.iii}
If $X=\Spec(A)$, where $A$ is a Noetherian ring, and
$\sh F=\widetilde M$, where $M$ is a finite $A$-module, we say that $M$
has property $(S_k)$ if $\sh F$ has this property.
For an arbitrary locally Noetherian prescheme $X$, to say that
$\sh F$ has property $(S_k)$ \emph{at a point} $x\in X$ therefore means
that, if we set $Y=\Spec(\sh O_x)$, the $\sh O_Y$-Module
$\widetilde{\mathscr F}_x$ has property $(S_k)$; note that condition
\eqref{IV.5.7.2.1} in general \emph{does not imply}
\eqref{IV.5.7.2.2} for every generization $x'$ of $x$, since there is
no relation of inequality between
$\operatorname{prof}(\sh F_x)$ and
$\operatorname{prof}(\sh F_{x'})$ \sref[0]{0.16.4.6}.

\item
\phantomsection
\label{IV.5.7.3.iv}
It follows immediately from the definition that if $\sh F$
satisfies $(S_k)$ at a point $x$, it also satisfies $(S_k)$ at every point
$x'$ that is a \emph{generization} of $x$.

\item
\phantomsection
\label{IV.5.7.3.v}
Property $(S_k)$ is especially important for $k=1$ and $k=2$; it was
introduced for $k=2$ by Serre, to express his normality
criterion; see \sref{IV.5.8.6}.\footnote{The printed French text points to
\textup{(5.8.5)}; Serre's normality criterion is Theorem
\textup{(5.8.6)}.}

\item
\phantomsection
\label{IV.5.7.3.vi}
Let $X$ be a locally Noetherian prescheme, $Y$ a closed sub-prescheme
of $X$, $j:Y\to X$ the canonical injection, and $\sh G$ a coherent
$\sh O_Y$-Module.  It is clear that for every $x\in Y$, we have
\[
\dim(\sh G_x)=\dim((j_*(\sh G))_x),\qquad
\operatorname{prof}(\sh G_x)
 =\operatorname{prof}((j_*(\sh G))_x),
\]
whence $\operatorname{coprof}(\sh G_x)
=\operatorname{coprof}((j_*(\sh G))_x)$.
For $\sh G$ to satisfy $(S_k)$, it is necessary and sufficient that
$j_*(\sh G)$ satisfy $(S_k)$.
\end{enumerate}
\end{remarks}

\begin{proposition}[5.7.4]
\phantomsection
\label{IV.5.7.4}
Let $X$ be a locally Noetherian prescheme and let $\sh F$ be a
coherent $\sh O_X$-Module.  Put
$S=\operatorname{Supp}(\sh F)$, and for every $n\geq0$ let $Z_n$ be
the set of points $z\in X$ such that
$\operatorname{coprof}(\sh F_z)>n$, so that $Z_n\subset S$.
\begin{enumerate}[{\rm(i)}]
\item
\phantomsection
\label{IV.5.7.4.i}
For $\sh F$ to satisfy property $(S_k)$, it is necessary and sufficient
that, for every $n\geq0$, we have
\begin{equation}
\operatorname{codim}(Z_n,S)>n+k.
\tag{5.7.4.1}\label{IV.5.7.4.1}
\end{equation}

\item
\phantomsection
\label{IV.5.7.4.ii}
Suppose in addition that the $Z_n$ are closed in $X$.  Then, for
$\sh F$ to satisfy property $(S_k)$ at a point $x\in S$, it is necessary
and sufficient that, for every $n\geq0$, we have
\begin{equation}
\operatorname{codim}_x(Z_n,S)>n+k.
\tag{5.7.4.2}\label{IV.5.7.4.2}
\end{equation}
\end{enumerate}
\end{proposition}

\begin{proof}
\begin{enumerate}[{\rm(i)}]
\item
Indeed, by definition \sref{IV.5.1.3}, $\operatorname{codim}(Z_n,S)
=\inf_{z\in Z_n}(\dim(\sh O_{S,z}))$,
and the inequality \eqref{IV.5.7.4.1} therefore means
\sref{IV.5.1.12.2} that, for every $z\in X$ and every $n\geq0$, the
relation
\[
\dim(\sh F_z)-\operatorname{prof}(\sh F_z)\geq n+1
\]
implies the relation
\[
\dim(\sh F_z)\geq n+k+1.
\]
But if we set $a=\dim(\sh F_z)$,
$b=\operatorname{prof}(\sh F_z)$, we have $b\leq a$, and to say that for
every $n\geq0$ the relation $b\leq a-n-1$ implies $k\leq a-n-1$ is
equivalent to saying that $b\geq\inf(k,a)$, whence the proposition.

\oldpage[IV-2]{105}

\item
The reasoning is the same as in
{\hyperref[IV.5.7.4.i]{\textup{(i)}}}, except that we must restrict
ourselves to the $z\in X$ which are \emph{generizations} of $x$, and
take \sref{IV.5.1.3.2} into account.
\end{enumerate}
\end{proof}

\begin{proposition}[5.7.5]
\phantomsection
\label{IV.5.7.5}
Let $X$ be a locally Noetherian prescheme and let $\sh F$ be a
coherent $\sh O_X$-Module.  For $\sh F$ to satisfy $(S_k)$
\textup{(}$k\geq1$\textup{)}, it is necessary and sufficient that, for
every integer $r$ such that $0\leq r<k$, every open subset $U$ of
$X$, and every $(\sh F|U)$-regular sequence
$(f_i)_{1\leq i\leq r}$ of sections of $\sh O_X$ over $U$,
$(\sh F|U)/(\sum_{i=1}^{r}f_i(\sh F|U))$
have no embedded associated prime cycle.
\end{proposition}

\begin{proof}
Let us first prove the proposition for $k=1$; it can then also be stated
by saying that \emph{for $\sh F$ to satisfy $(S_1)$, it is necessary and
sufficient that $\sh F$ have no embedded associated prime cycle}.  Indeed,
to say that $\sh F$ satisfies $(S_1)$ means that at the points
$x\in\operatorname{Supp}(\sh F)$ such that $\dim(\sh F_x)>0$, we have
$\operatorname{prof}(\sh F_x)\geq1$ (since at the other points of
$\operatorname{Supp}(\sh F)$ we have $\dim(\sh F_x)=0$, hence
$\operatorname{prof}(\sh F_x)=0$); but to say that
$\operatorname{prof}(\sh F_x)\geq1$ means that $\mathfrak m_x$
\emph{is not associated with} $\sh F_x$ \sref[0]{0.16.4.6},~(i), or
again that $x$ is not associated with $\sh F$ \sref{IV.3.1.1}; on the
other hand, if $S$ is a sub-prescheme of $X$ having
$\operatorname{Supp}(\sh F)$ as underlying space, we have
\sref{IV.5.1.12.1} $\dim(\sh F_x)=\dim(\sh O_{S,x})$; to say that
$\dim(\sh F_x)>0$ therefore means that $x$ is not a maximal point of
$\operatorname{Supp}(\sh F)$ \sref{IV.5.1.2}, whence the conclusion.

Secondly, let us show that, for arbitrary $k>1$, the condition of the
statement is \emph{necessary}.  We may restrict ourselves to the case
where $U=X$, and our assertion will be proved (by induction on $k$ and by
virtue of the first part of the reasoning) if we show that when $\sh F$
satisfies $(S_k)$ ($k>1$) and $f$ is an $\sh F$-\emph{regular} section
of $\sh O_X$ over $X$, then $\sh F/f\sh F$ \emph{satisfies}
$(S_{k-1})$ (i.e.\ satisfies \eqref{IV.5.7.2.1} with $k$ replaced by
$k-1$).  Now, for every $x\in X$, $f_x$ is an $\sh F_x$-regular element;
if it is invertible, we have $\sh F_x/f_x\sh F_x=0$ and the conclusion is
trivial.  If on the contrary $f_x\in\mathfrak m_x$, we know that
$\operatorname{prof}(\sh F_x/f_x\sh F_x)
=\operatorname{prof}(\sh F_x)-1$ \sref[0]{0.16.4.6},~(i), and
$\dim(\sh F_x/f_x\sh F_x)=\dim(\sh F_x)-1$ \sref[0]{0.16.3.4}, and our
assertion follows.

Finally, let us prove that for $k>1$, the condition of the statement is
\emph{sufficient}.  We shall proceed by induction on $k$; let $x$ be a
point of $X$, and suppose first that $\dim(\sh F_x)\geq k$.  The
induction hypothesis implies that $\sh F$ satisfies $(S_{k-1})$, hence
$\operatorname{prof}(\sh F_x)\geq k-1$; taking \sref[0]{0.15.2.4} into
account, there is therefore an open neighbourhood $U$ of $x$ and an
$(\sh F|U)$-regular sequence $(f_i)_{1\leq i\leq k-1}$ of sections of
$\sh O_X$ over $U$ such that $(f_i)_x\in\mathfrak m_x$ for
$1\leq i\leq k-1$.  The hypothesis implies that
$\sh G=(\sh F|U)/(\sum_{i=1}^{k-1}f_i(\sh F|U))$
has no embedded associated prime cycle; but we have
$\dim(\sh G_x)=\dim(\sh F_x)-(k-1)\geq1$ \sref[0]{0.16.3.4}, so $x$
is not associated with $\sh G$; we therefore have
$\operatorname{prof}(\sh G_x)\geq1$ \sref[0]{0.16.4.6}, and since
$\operatorname{prof}(\sh G_x)=\operatorname{prof}(\sh F_x)-(k-1)$
\sref[0]{0.16.4.6}, we have $\operatorname{prof}(\sh F_x)\geq k$.
Suppose secondly that $\dim(\sh F_x)=r<k$; since $\sh F$ satisfies
$(S_{k-1})$, we have
$\operatorname{prof}(\sh F_x)\geq\inf(k-1,r)=r$, and this completes the
proof.
\end{proof}

\begin{corollary}[5.7.6]
\phantomsection
\label{IV.5.7.6}
Suppose that $\sh F$ satisfies $(S_k)$, where $k\geq1$.  If
$(f_i)_{1\leq i\leq r}$ is an $\sh F$-regular sequence of sections
of $\sh O_X$ over $X$ and $r<k$, then $\sh F\Big/\Big(\sum_{i=1}^{r}f_i\sh F\Big)$
satisfies $(S_{k-r})$.
\end{corollary}

\begin{proof}
This follows immediately from \sref{IV.5.7.5}.
\end{proof}

\begin{corollary}[5.7.7]
\phantomsection
\label{IV.5.7.7}
For $\sh F$ to satisfy $(S_2)$, it is necessary and sufficient that
$\sh F$ have no embedded associated prime cycle and that, for every open
subset $U$ of $X$ and every $(\sh F|U)$-regular section $f$ of $\sh O_X$
over $U$, $(\sh F|U)/f(\sh F|U)$ have no embedded associated prime
cycle.
\end{corollary}

\oldpage[IV-2]{106}

\begin{remarks}[5.7.8]
\phantomsection
\label{IV.5.7.8}
Let $X$ be a locally Noetherian prescheme of dimension $1$.  The
following conditions are equivalent: $X$ is Cohen--Macaulay; $X$
satisfies $(S_1)$; $X$ satisfies $(S_n)$ for some $n\geq1$.  This
follows from Definitions \sref{IV.5.7.1} and \sref{IV.5.7.2}.
By \sref{IV.5.7.5}, these conditions are also equivalent to the
absence of an embedded associated prime cycle.  In particular,
every reduced locally Noetherian prescheme of dimension $1$ is
Cohen--Macaulay.
\end{remarks}

\begin{proposition}[5.7.9]
\phantomsection
\label{IV.5.7.9}
Let $A$ and $B$ be Noetherian rings, let $\rho:A\to B$ be a ring
homomorphism, and let $M$ be a $B$-module such that the
$A$-module $M_{[\rho]}$ is finite.  Let $\mathfrak p$ be a prime
ideal of $A$.  The prime ideals of $B$ lying over $\mathfrak p$ and
belonging to $\operatorname{Supp}(M)$ are finite in number.  If
$(\mathfrak q_i)_{1\leq i\leq n}$ is the family of these ideals,
then
\begin{align}
\dim_{A_{\mathfrak p}}(M_{\mathfrak p})
 &=\sup_i\dim_{B_{\mathfrak q_i}}(M_{\mathfrak q_i})
\tag{5.7.9.1}\label{IV.5.7.9.1}\\
\operatorname{prof}_{A_{\mathfrak p}}(M_{\mathfrak p})
 &=\inf_i\operatorname{prof}_{B_{\mathfrak q_i}}
   (M_{\mathfrak q_i}).
\tag{5.7.9.2}\label{IV.5.7.9.2}
\end{align}
\end{proposition}

\begin{proof}
Let $\mathfrak b$ be the annihilator of $M$.  The ring
$B/\mathfrak b$ identifies with an $A$-submodule of
$\operatorname{End}_A(M_{[\rho]})$, and is therefore finite over
$A$.  Hence only finitely many prime ideals of $B$ lie over
$\mathfrak p$ and contain $\mathfrak b$; these are precisely the
ones belonging to $\operatorname{Supp}(M)$.

Replacing $B$ by $B/\mathfrak b$ does not change the right-hand
sides of \eqref{IV.5.7.9.1} and \eqref{IV.5.7.9.2}
\sref[0]{0.16.1.9}, \sref[0]{0.16.4.8}; we may therefore suppose
that $B$ is finite over $A$.  Put $S=A\setmin\mathfrak p$ and
$B'=S^{-1}B$.  Then $B'$ is a Noetherian semilocal ring whose
maximal ideals are $S^{-1}\mathfrak q_i$, and $\dim_{A_{\mathfrak p}}(M_{\mathfrak p})
=\dim_{B'}(M_{\mathfrak p})$
\sref[0]{0.16.1.9}.  Formula \eqref{IV.5.7.9.1} is now the
semilocal dimension formula \sref[0]{0.16.1.7.4}.

For \eqref{IV.5.7.9.2}, as in \sref[0]{0.16.4.8}, it is enough to
treat the case
$\operatorname{prof}_{A_{\mathfrak p}}(M_{\mathfrak p})=0$.  The
same argument shows that $M_{\mathfrak p}$ contains a
$B'$-submodule of finite length, and hence a simple
$B'$-submodule.  Such a module is isomorphic to the residue field
of one of the local rings $B_{\mathfrak q_i}$.  Thus $\operatorname{prof}_{B_{\mathfrak q_i}}(M_{\mathfrak q_i})=0$
for at least one $i$ \sref[0]{0.16.4.6}, proving the formula.
\end{proof}

\begin{corollary}[5.7.10]
\phantomsection
\label{IV.5.7.10}
Assume the hypotheses of \sref{IV.5.7.9}, and suppose in addition
that $A$ is \emph{local}.  The $A$-module $M_{[\rho]}$ is Cohen--Macaulay
if and only if every $M_{\mathfrak q_i}$, for
$\mathfrak q_i$ lying over the maximal ideal of $A$, is a
Cohen--Macaulay $B_{\mathfrak q_i}$-module and all the numbers
$\dim_{B_{\mathfrak q_i}}(M_{\mathfrak q_i})$ are equal.
\end{corollary}

\begin{proof}
By \eqref{IV.5.7.9.1} and \eqref{IV.5.7.9.2}, these conditions are
equivalent to
\[
\dim_A(M_{[\rho]})
 =\operatorname{prof}_A(M_{[\rho]}).
\]
\end{proof}

\begin{corollary}[5.7.11]
\phantomsection
\label{IV.5.7.11}
Assume the hypotheses of \sref{IV.5.7.9}.
\begin{enumerate}[{\rm(i)}]
\item
\phantomsection
\label{IV.5.7.11.i}
If $M_{[\rho]}$ satisfies $(S_k)$, then $M$ satisfies $(S_k)$.

\item
\phantomsection
\label{IV.5.7.11.ii}
Suppose that for every pair of prime ideals
$\mathfrak q,\mathfrak q'$ of $B$ satisfying
$\rho^{-1}(\mathfrak q)=\rho^{-1}(\mathfrak q')$, one has $\dim_{B_{\mathfrak q}}(M_{\mathfrak q})
=\dim_{B_{\mathfrak q'}}(M_{\mathfrak q'})$.
If $M$ satisfies $(S_k)$, then $M_{[\rho]}$ satisfies $(S_k)$.
\end{enumerate}
\end{corollary}

\oldpage[IV-2]{107}

\begin{proof}
This follows immediately from \eqref{IV.5.7.9.1},
\eqref{IV.5.7.9.2}, and the definition of $(S_k)$.
\end{proof}

\begin{env}[5.7.12]
\phantomsection
\label{IV.5.7.12}
In accordance with the definitions of \sref{IV.5.7.1}, given an
\emph{arbitrary} Noetherian ring $A$ and a finite $A$-module $M$, we define
\[
\operatorname{coprof}_A(M)
 =\operatorname{coprof}(M)
 =\sup_{x\in X}(\operatorname{coprof}_{A_x}(M_x)),
\qquad X=\Spec(A).
\]
We shall see later \sref{IV.6.11.5} that this definition \emph{agrees} with
that of \sref[0]{0.16.4.9} when $A$ is a Noetherian \emph{local} ring.
\end{env}

\begin{corollary}[5.7.13]
\phantomsection
\label{IV.5.7.13}
Let $A$ and $B$ be Noetherian rings, let $\rho:A\to B$ be a ring
homomorphism, and let $M$ be a $B$-module such that $M_{[\rho]}$ is
a finite $A$-module.  Then
\begin{equation}
\operatorname{coprof}_A(M_{[\rho]})
 \geq \operatorname{coprof}_B(M).
\tag{5.7.13.1}\label{IV.5.7.13.1}
\end{equation}
\end{corollary}

\begin{proof}
This follows from \sref{IV.5.7.12} and
\eqref{IV.5.7.9.1}, \eqref{IV.5.7.9.2}.
\end{proof}
